\title{FlashAttention-3: Fast and Accurate Attention with Asynchrony and Low-precision}

\begin{abstract}
The attention mechanism forms the foundational component within Transformer architectures, but also presents a key computational bottleneck, particularly for large language models and applications that involve extended contexts. \textsc{FlashAttention} introduced an optimized scheme to accelerate attention computations on GPUs by reducing the frequency of memory access operations. Nevertheless, this approach does not fully utilize the advanced functionalities of cutting-edge hardware, as evidenced by \textsc{FlashAttention-2}, which operates at merely 35\% efficiency on H100 GPUs. In this work, we propose three principal strategies for enhancing attention performance on Hopper GPUs: leveraging asynchrony between the Tensor Cores and TMA to (1) synchronize computation with data transfer through warp-specialization and (2) execute block-wise matmul and softmax in an interleaved manner; and (3) employing block quantization and incoherent processing to exploit the hardware's native support for low-precision FP8 format. Our proposed method, \textsc{FlashAttention-3}, delivers a performance boost of 1.5-2.0$\times$ on H100 GPUs—achieving up to 740 TFLOPs/s using FP16 (representing 75\% utilization)—and approaches 1.2 PFLOPs/s in FP8 precision. Moreover, empirical analysis shows that FP8 \textsc{FlashAttention-3} offers a 2.6$\times$ reduction in numerical error relative to a standard FP8 attention baseline.
\end{abstract}

\section{Introduction}
\label{sec:intro}

Within the Transformer framework~\citep{vaswani2017attention}, the main computational constraint arises from the attention mechanism, as the calculation of self-attention scores across queries and keys grows quadratically with respect to the sequence length. Enhancing attention mechanisms to accommodate longer contexts promises to enable advancements such as improved modeling and inference over extended textual content~\citep{guo2021longt5,shaham2022scrolls,peng2023yarn} and handling multiple files in expansive code repositories~\citep{roziere2023code, li2023starcoder}, as well as facilitating novel modalities, including high-resolution imagery~\citep{chen2022scaling}, audio streams~\citep{gulati2020conformer}, and video data~\citep{ho2022video}. Furthermore, this progress extends to innovative applications, such as systems managing lengthy user interaction histories~\citep{sun2019bert4rec} and agents capable of executing tasks with extended temporal dependencies~\citep{yao2022react}. These opportunities have spurred extensive research into accelerating attention in settings with long input sequences, via methods of approximation~\citep{katharopoulos2020transformers,choromanski2020rethinking, tay2020efficient}, improvements in software implementation~\citep{rabe2021self,dao2022flashattention,kwon2023efficient}, and exploration of fundamentally different architectures~\citep{peng2023rwkv,sun2023retentive,gu2023mamba}.

This study extends the contributions of \citet{dao2022flashattention} in formulating exact-attention methods that explicitly leverage the architecture and operational nuances of GPUs within their design framework. Dao et al., in \citep{dao2022flashattention}, proposed \textsc{FlashAttention}, which employs an innovative tiling approach to parallelize attention by combining all associated operations into a single GPU kernel, thereby minimizing unnecessary accesses to slower global memory for intermediate data.

\citet{dao2023flashattention2} further evolved the algorithm with \textsc{FlashAttention-2}, which introduces parallelization across the sequence length and reconfigures the forward pass to process blocks of the key and value matrices, thereby increasing workload distribution and optimizing GPU occupancy. Nevertheless, we find that \textsc{FlashAttention-2} does not reach optimal resource utilization on the latest GPUs compared to finely-tuned GEMM kernels—utilization levels are often 35\% versus the 80-90\% observed on Hopper H100 for GEMM. This discrepancy may be partially due to implementation-related issues, including reliance on Ampere-era instructions rather than Hopper-specific ones when utilizing Tensor Cores. Notably, recent advancements such as ThunkerKitten~\citep{spector2024thunder} and cuDNN 9~\citep{cudnn9} indicate that applying Hopper-oriented instructions and tile abstractions can accelerate attention processing and streamline the development process.

At a foundational level, the algorithm underlying \textsc{FlashAttention-2} is constructed around a straightforward synchronous paradigm and does not intentionally leverage asynchrony or low-precision computation as part of its design philosophy. The emergence of asynchrony is attributable to hardware advancements, where specialized units—such as Tensor Cores for matrix multiplication and Tensor Memory Accelerator (TMA) for data transfer—operate independently from conventional CUDA cores responsible for logic and various arithmetic operations. Meanwhile, adoption of lower numerical precision formats, including FP8 on Hopper and FP4 on Blackwell, as successors to FP16 (introduced with Pascal in 2017) and BF16 (on Ampere in 2020), has repeatedly demonstrated substantial throughput gains, doubling or quadrupling computation within the same energy and silicon footprint. In \cref{subsec:hardware}, we discuss in detail how Hopper supports these performance enhancements.

The main technical obstacle is refactoring \textsc{FlashAttention-2} to exploit these capabilities: embracing asynchrony necessitates concurrent execution of matrix multiplication and softmax operations—even when dependencies exist between them—while employing low-precision computation demands careful strategies to contain quantization errors, which are particularly pronounced with outlier feature values in LLMs~\citep{dettmers2208llm, sun2024massive}.

In pursuit of this goal, we introduce \textsc{FlashAttention-3}, which integrates and advances three novel strategies to enhance efficiency on the latest GPU platforms.\footnote{While our findings are discussed with reference to NVIDIA's Hopper architecture, the proposed algorithm is applicable to any GPU system that supports highly capable asynchronous operations and low-precision processing.}

\begin{enumerate}

\item \textbf{Producer-Consumer asynchrony:} Our approach employs warp-specialized software pipelining, which capitalizes on the asynchronous nature of data transfer and Tensor Core execution by partitioning data producers and consumers across distinct warps. This design enhances the algorithm's capability to mask both memory access delays and instruction scheduling latency.
\item \textbf{Hiding softmax under asynchronous block-wise GEMMs:} To increase efficiency, we schedule non-GEMM computations associated with softmax—such as floating-point multiply-add and exponentiation—to occur in parallel with asynchronous WGMMA instructions that carry out GEMMs. The \textsc{FlashAttention-2} algorithm is modified to eliminate certain serial dependencies between the softmax calculations and GEMM operations. For instance, in the 2-stage variant of the method, softmax is performed on one section of the scores matrix concurrently as WGMMA asynchronously computes the subsequent section.
\item \textbf{Hardware-accelerated low-precision GEMM:} The forward pass routine is redesigned to enable the use of FP8 Tensor Cores in GEMM operations, resulting in nearly a twofold increase in observed TFLOPs/s. Addressing the disparate memory layout constraints of WGMMA for blocks of FP32 accumulators and FP8 operands is essential, and this is handled through block quantization and incoherent processing techniques, which help compensate for the accuracy reduction inherent to FP8 arithmetic.
\end{enumerate}

To assess the empirical performance of our approach, we evaluate \textsc{FlashAttention-3} on the H100 SXM5 GPU across various settings and report the following findings: (1) The FP16 configuration delivers a speed increase of 1.5--2.0$\times$ over \textsc{FlashAttention-2} for the forward pass (peaking at 740 TFLOPs/s), and 1.5--1.75$\times$ during the backward pass, (2) FP8 achieves nearly 1.2 PFLOPs/s throughput, and (3) for extended sequence lengths, FP16 excels while FP8 remains competitive\footnote{Specifically, at head dimension 64, FP8 in \textsc{FlashAttention-3} surpasses alternatives; for head dimensions 128 and 256, performance matches cuDNN's attention in the absence of causal masking, but trails it with causal masking.} when compared with the cutting-edge attention implementation from NVIDIA's cuDNN library. Our results further confirm that FP16 \textsc{FlashAttention-3} maintains the same numerical accuracy as \textsc{FlashAttention-2}, and offers superior error characteristics relative to conventional attention, as intermediate computations (such as softmax normalization) remain in FP32. Additionally, FP8 \textsc{FlashAttention-3} equipped with block quantization and incoherent processing achieves 2.6$\times$ higher accuracy than standard attention utilizing per-tensor quantization in scenarios involving outlier features.

We release \textsc{FlashAttention-3} under a flexible license\footnote{\textsc{FlashAttention-3}
  is available at \texttt{https://github.com/Dao-AILab/flash-attention}}, and are committed to incorporating it into both PyTorch and Hugging Face frameworks, thereby maximizing its accessibility and utility for the research and development communities.

\section{Background: Multi-Head Attention and GPU Characteristics}
\label{sec:background}

\subsection{Multi-Head Attention}
\label{subsec:multi_head_attn}

Let $\mathbf{Q}, \mathbf{K}, \mathbf{V} \in \mathbb{R}^{N \times d}$ be the query, key and value input sequences associated to a single head, where $N$ is the sequence length and $d$ is the head dimension. Then the attention output $\mathbf{O}$ is computed as:

\begin{equation*}
  \mathbf{S} = \alpha \mathbf{Q} \mathbf{K}^\top \in \mathbb{R}^{N \times N}, \quad \mathbf{P} = \mathrm{softmax}(\mathbf{S}) \in \mathbb{R}^{N \times N}, \quad \mathbf{O} = \mathbf{P}\mathbf{V} \in \mathbb{R}^{N \times d},
\end{equation*}

The $\mathrm{softmax}$ operation is performed along each row, and the scaling factor $\alpha$ is generally chosen to be $1/\sqrt{d}$. To ensure stable numerical computations when applying the exponential, $\mathrm{rowmax}(\mathbf{S})$ is subtracted from $\mathbf{S}$. In the case of multi-head attention (MHA), every attention head utilizes distinct projections for queries, keys, and values. This allows the calculations to be carried out in parallel for different heads and batches, ultimately generating the complete output tensor.

Now let $\phi$ be a scalar loss function and let $\mathbf{d}(-) = \partial \phi / \partial (-)$ be notation for the gradient. Given the output gradient $\mathbf{dO} \in \mathbb{R}^{N \times d}$, we compute $\mathbf{dQ}$, $\mathbf{dK}$, and $\mathbf{dV}$ according to the chain rule as follows:

\begin{align*}
  \mathbf{dV} &= \mathbf{P}^\top \mathbf{dO} \in \mathbb{R}^{N \times d} \\
  \mathbf{dP} &= \mathbf{dO} \mathbf{V}^\top \in \mathbb{R}^{N \times N} \\
  \mathbf{dS} &= \mathrm{dsoftmax} (\mathbf{dP}) \in \mathbb{R}^{N \times N} \\
  \mathbf{dQ} &= \alpha \mathbf{dS} \mathbf{K} \in \mathbb{R}^{N \times d} \\
  \mathbf{dK} &= \alpha \mathbf{dS}^\top \mathbf{Q} \in \mathbb{R}^{N \times d},
\end{align*}

In this context, the relationship $\mathbf{d}s = (\mathrm{diag}(p) - p p^\top)\mathbf{d}p$ holds, where $p = \mathrm{softmax}(s)$ is viewed as a function of the vector $s$. We denote by $\mathrm{dsoftmax}(\mathbf{dP})$ the application of this operation to each row individually. Importantly, for the backward propagation step in MHA, this procedure is likewise efficiently parallelizable across all heads and batches.

\subsection{GPU hardware characteristics and execution model}
\label{subsec:hardware}

We discuss elements of the GPU execution model pertinent to \textsc{FlashAttention-3}, concentrating specifically on the NVIDIA Hopper architecture to exemplify these concepts.

\paragraph{Memory hierarchy:} GPU memories are structured into a hierarchical system of data storage locations, where higher bandwidth typically corresponds to smaller capacities (\cref{tab:gpu-hierarchy})\footnote{\citet{luo2024benchmarking} identifies shared memory bandwidth as 128 bytes per clock cycle per SM; for total bandwidth, this is scaled by the presence of 132 SMs and a boost clock speed of 1830 MHz.}. The off-chip DRAM, referred to as global memory (GMEM) or HBM, is universally accessible by all streaming multiprocessors (SMs). Data transferred from GMEM is automatically placed into the on-chip L2 cache. Within each SM, there is a limited-capacity, on-chip cache—shared memory (SMEM)—which features extensive banking and explicit control by the programmer. The innermost tier consists of the register file that resides inside every SM.

\paragraph{Thread hierarchy:} The structure of the GPU programming model is centered on the logical arrangement of computational units known as threads. This hierarchy, moving from the smallest to the largest grouping, consists of threads, warps (each containing 32 threads), warpgroups (formed by 4 neighboring warps), threadblocks (also referred to as cooperative thread arrays or CTAs), threadblock clusters (present in Hopper architecture), and grids at the topmost level.

The interconnection between these hierarchies is strong. Threads belonging to a particular CTA are concurrently executed on the same SM, while CTAs within a shared cluster are scheduled together on a single GPC. All threads within a CTA can directly access SMEM, while individual threads have exclusive access to a maximum of 256 registers (RMEM).

\paragraph{Asynchrony and warp-specialization:}

GPUs function as throughput-oriented processors, utilizing parallelism and asynchronous operations to mitigate latency arising from memory accesses and instruction execution. To facilitate asynchronous data transfers between GMEM and SMEM, Hopper integrates the Tensor Memory Accelerator (TMA), a specialized hardware component \cite[\S7.29]{cuda}. Moreover, Hopper’s Tensor Core, accessed through the warpgroup-level WGMMA instruction \cite[\S9.7.14]{ptx}, distinguishes itself from previous designs such as Ampere by providing asynchronous capabilities and permitting direct input reads from shared memory.

With hardware mechanisms for asynchrony, it becomes possible to implement warp-specialized kernels in which the warps within a CTA are assigned dedicated producer or consumer responsibilities, focusing exclusively on either data transfer or computation tasks. This general approach enhances the compiler’s capacity to devise more effective instruction schedules \citep{warp-specialization-2011}. Furthermore, Hopper introduces the ability to dynamically redistribute registers among warpgroups by means of \verb|setmaxnreg| \cite[\S9.7.17.1]{ptx}, enabling warps that perform MMAs to access a greater portion of RMEM compared to warps dedicated to TMA, which typically require just a single thread.

\paragraph{Low-precision number formats:}
\label{sec:low-precision-gpu}
Contemporary GPUs incorporate dedicated hardware modules designed to optimize computations involving low-precision formats. For instance, on Hopper architecture, utilizing the WGMMA instruction with FP8 Tensor Cores achieves double the throughput per SM relative to FP16 or BF16 precision.

To effectively utilize FP8 WGMMA, it is essential to be aware of the required operand layout constraints. In the case of a GEMM operation involving $A \times B^{\top}$, where $A$ is an $M\times K$ matrix and $B$ is $N\times K$, the term \emph{mn-major} refers to an operand with contiguous memory layout along the $M$ or $N$ dimension, while \emph{k-major} indicates contiguity in the $K$ dimension. With FP16 WGMMA, operands stored in SMEM may be presented in either mn-major or k-major format. In contrast, FP8 WGMMA restricts acceptable input operands to exclusively the k-major configuration. This constraint becomes particularly problematic in scenarios—such as attention mechanisms—that require chaining multiple GEMMs within a single kernel; mismatched FP32 accumulator layouts and FP8 operand memory arrangements hinder direct invocation of successive, dependent FP8 WGMMAs.

Regarding attention mechanisms, these layout constraints necessitate specific adjustments to the structure of an FP8 algorithm. Details of these adaptations are presented in \cref{sec:algofp8}.

\subsection{Standard Attention and Flash Attention} As outlined by \citet{dao2022flashattention}, the term \textbf{standard attention} refers to a GPU-based attention computation that explicitly stores the intermediate matrices $\mathbf{S}$ and $\mathbf{P}$ in HBM. The core innovation of \textsc{FlashAttention} involves employing a local variant of softmax reduction, thereby eliminating costly intermediate memory accesses and integrating the attention computation within a single fused kernel. This local softmax process is realized in lines \ref{code-ws:softmax_start}-\ref{code-ws:softmax_end} of the consumer mainloop from \cref{alg:flash3_wgmma_ws_only}, combined with block-wise rescaling of $\mathbf{O}$. A concise proof that this method yields $\mathbf{O}$ is provided in \cite[\S 2.3.1]{dao2023flashattention2}.

\section{FlashAttention-3: Algorithm}
\label{sec:algo}

This section provides an overview of the \textsc{FlashAttention-3} algorithm. We concentrate on the forward computation here, with details regarding the backward pass presented in~\cref{sec:algo_ws_bwd}. Initially, we outline the integration of warp-specialization and a circular SMEM buffer into the foundational framework of \textsc{FlashAttention-2}. Subsequently, we demonstrate how to leverage WGMMA's asynchrony to construct a two-stage pipeline that overlaps GEMM and softmax operations. Lastly, we delineate the necessary adjustments for FP8, addressing both layout requirements and accuracy through block quantization and incoherent processing.

\subsection{Producer-Consumer asynchrony through warp-specialization and pingpong scheduling}
\label{sec:algo_ws}

\paragraph{Warp-specialization}
Similar to the approach taken in \textsc{FlashAttention-2}, the forward computation in \textsc{FlashAttention-3} is highly parallelizable along dimensions such as batch size, head count, and the length of the query sequence. Consequently, it is sufficient to describe the algorithm at the CTA level, where the core operation focuses on processing a tile $\mathbf{Q}_i$ of the query matrix in order to produce the corresponding tile $\mathbf{O}_i$ for the output. For clarity, the following presentation outlines the warp-specialization methodology assuming a circular SMEM buffer is used, \emph{excluding} any overlapping between the GEMM and softmax steps. Define $d$ as the head dimension, $N$ as the sequence length, and select a query block size $B_r$ so that the query matrix $\mathbf{Q}$ can be partitioned into $T_r = \lceil \frac{N}{B_r} \rceil$ blocks, $\mathbf{Q}_1, .., \mathbf{Q}_{T_r}$.

\begin{algorithm}[H]
    \caption{\small\label{alg:flash3_wgmma_ws_only}\textsc{FlashAttention-3} forward pass \textbf{without} intra-consumer overlapping -- CTA view}
    \begin{algorithmic}[1]
\REQUIRE Let $\mathbf{Q}_i \in \mathbb{R}^{B_r \times d}$ and $\mathbf{K}, \mathbf{V} \in \mathbb{R}^{N \times d}$ reside in HBM; let the key block have size $B_c$, and denote $T_c = \lceil \frac{N}{B_c} \rceil$.
\STATE Create a pipeline controller responsible for barrier synchronization, utilizing a $s$-stage circular SMEM buffer.
\IF {operating in producer warpgroup}
\STATE Free a predetermined set of registers.
\STATE Trigger a data transfer of $\mathbf{Q}_i$ from HBM to shared memory.
\STATE Once data is available, send a completion signal to the consumers indicating that $\mathbf{Q}_i$ is loaded.
\FOR{$0 \le j < T_c$}
    \STATE Stall until the $(j\,\%\,s)$th buffer stage has been processed by the consumers.
    \STATE Transfer $\mathbf{K}_j, \mathbf{V}_j$ from HBM to the shared memory segment corresponding to buffer stage $(j\,\%\,s)$.
    \STATE When load finishes, signal consumers that $\mathbf{K}_j, \mathbf{V}_j$ are now available.
\ENDFOR
\ELSE \STATE Reserve registers in an amount depending on how many consumer warps are active.
\STATE Allocate on-chip buffers: set $\mathbf{O}_i = (0) \in \mathbb{R}^{B_r \times d}$; also set $\ell_i = 0$ and $m_i = -\infty$ in $\mathbb{R}^{B_r}$.
\STATE Wait until $\mathbf{Q}_i$ appears in shared memory.
\FOR{$0 \le j < T_c$}
\STATE Wait until shared memory load of $\mathbf{K}_j$ completes.
\STATE Compute $\mathbf{S}_i^{(j)} = \mathbf{Q}_i \mathbf{K}_j^T$ by SS-GEMM. Upon finishing, synchronize. \label{alg:ws_only_gemm1}
\STATE Cache $m_i^{\mathrm{old}} = m_i$ and update $m_i \leftarrow \mathrm{max}(m_i^{\mathrm{old}}, \mathrm{rowmax}(\mathbf{S}_i^{(j)}))$. \label{code-ws:softmax_start}
\STATE Obtain $\widetilde{\mathbf{P}}_i^{(j)} = \mathrm{exp}(\mathbf{S}_i^{(j)} - m_i)$; update $\ell_i = \mathrm{exp}(m_i^{\mathrm{old}} - m_i) \ell_i + \mathrm{rowsum}(\widetilde{\mathbf{P}}_i^{(j)})$. \label{code-ws:softmax_end}
\STATE Wait until shared memory provides $\mathbf{V}_j$.
\STATE Update $\mathbf{O}_i = \mathrm{diag}(\mathrm{exp}(m_i^{\mathrm{old}} - m_i))^{-1} \mathbf{O}_i + \widetilde{\mathbf{P}}_i^{(j)} \mathbf{V}_j$ using RS-GEMM. Synchronize on completion. \label{alg:ws_only_gemm2}
\STATE Free the buffer stage $(j\,\%\,s)$ so that the producer can proceed.
\ENDFOR
\STATE Normalize $\mathbf{O}_i$ with $\mathbf{O}_i = \mathrm{diag}(\ell_i)^{-1} \mathbf{O}_i$; set $L_i = m_i + \log(\ell_i)$.
\STATE Store the $i$th segments $\mathbf{O}_i$ and $L_i$ into the global memory locations for $\mathbf{O}$ and $L$.
\ENDIF
\end{algorithmic}
\end{algorithm}

In our deployment of \cref{alg:flash3_wgmma_ws_only} on Hopper, memory (de)allocation is managed through \verb|setmaxnreg|, while TMA is utilized for loading $\mathbf{Q}_i$ as well as $\{ \mathbf{K}_j, \mathbf{V}_j \}_{0 \leq j < T_c}$. WGMMA carries out the GEMM operations within the consumer mainloop, with the SS or RS prefixes denoting whether the initial operand resides in shared memory or in the register file, respectively. To facilitate understanding of the execution pattern of \cref{alg:flash3_wgmma_ws_only}, it is important to recognize that TMA load instructions do not block each other due to their asynchronous nature. Additionally, during the producer mainloop, there will be no wait instructions for the first $s$ iterations, since these iterations correspond to the process of filling the buffer.

\paragraph{Pingpong scheduling} The asynchronous execution facilitated by WGMMA and TMA, combined with warp-specialization, creates possibilities for concurrently performing the softmax calculation in one warpgroup while carrying out GEMM operations in another. This concurrent scheduling is especially relevant due to the disparity in throughput between matmul and non-matmul operations on contemporary hardware accelerators. For instance, the H100 SXM5 GPU is capable of achieving 989 TFLOPS in FP16 matmul, whereas special functions such as the exponential function\footnote{According to the CUDA programming guide, a streaming multiprocessor (SM) can execute 16 special function operations per clock cycle. Using 16 operations, 132 SMs, and a 1830 MHz clock frequency results in a special function throughput of 3.9 TFLOPS.} (required for softmax) are limited to 3.9 TFLOPS. In an attention forward pass using FP16 and a head dimension of 128, the number of matmul FLOPS is 512 times greater than the exponential operations, yet the exponential throughput is 256 times lower, leading to scenarios where exponential calculations may occupy up to 50\% of the runtime relative to matmul. This imbalance becomes more pronounced in FP8 configurations, where matmul throughput doubles but the exponential throughput remains constant.

Because the exponential operation is carried out by a dedicated hardware component (the multi-function unit), it is most efficient to time the exponential computation to overlap with the matmul operations handled by the Tensor Cores. To achieve this overlap, we insert synchronization barriers (\texttt{bar.sync} instructions) to ensure that the GEMMs (specifically, GEMM1 -- $\mathbf{P} \mathbf{V}$ for a given iteration, and GEMM0 -- $\mathbf{Q} \mathbf{K}^\top$
for the subsequent iteration) belonging to warpgroup 1 are executed ahead of the GEMMs associated with warpgroup 2. Consequently, this arrangement allows the softmax computation in warpgroup 1 to proceed concurrently with GEMM operations in warpgroup 2. The process then alternates, such that warpgroup 2 proceeds with the softmax while warpgroup 1 resumes GEMMs, effectively realizing a ``pingpong'' scheduling approach. This scheduling method is depicted in~\cref{fig:pingpong_scheduling}. Although in reality the transitions between roles are not as seamless as the illustration suggests, this scheduling technique typically results in better performance—e.g., increasing throughput from 570 TFLOPS to 620–640 TFLOPS in FP16 forward passes with a head dimension of 128 and sequence length of 8192.

\paragraph{Attention variants} When implementing multi-query attention~\citep{shazeer2019fast} and grouped query attention~\citep{ainslie2023gqa}, we employ the method used in \textsc{FlashAttention-2}, modifying tensor indexing so as to prevent unnecessary replication of $\mathbf{K}$ and $\mathbf{V}$ within HBM.

\subsection{Intra-warpgroup overlapping GEMMs and softmax}

It is possible to achieve partial overlap of certain instructions between the softmax and the GEMMs, even inside a single warpgroup. Here, we present a specific method to accomplish this.

In the attention algorithm, operations within the inner loop (main loop) have sequential dependencies that impede parallelization within a single iteration. For example, (local) softmax (lines \ref{code-ws:softmax_start} to \ref{code-ws:softmax_end}) relies on the output $\mathbf{S}_i^{(j)}$ of the first GEMM, while the second GEMM takes its result $\widetilde{\mathbf{P}}_i^{(j)}$ as an operand. Indeed, the wait statements in lines \ref{alg:ws_only_gemm1} and \ref{alg:ws_only_gemm2} of \cref{alg:flash3_wgmma_ws_only} serialize the execution of softmax and GEMMs. However, we can break these dependencies by pipelining across iterations through additional buffers in registers. Pursuing this idea, we propose the following two-stage\footnote{Note that the number of stages of the overlapping scheme is bounded by, but need not equal, the number $s$ of stages in the circular SMEM buffer.} GEMM-softmax pipelining algorithm:

\begin{algorithm}[H]
    \caption{\small\label{alg:flash3_wgmma}\textsc{FlashAttention-3} consumer warpgroup forward pass}
    \begin{algorithmic}[1]
      \REQUIRE  Matrices $\mathbf{Q}_i \in \mathbb{R}^{B_r \times d}$ and $\mathbf{K}, \mathbf{V} \in \mathbb{R}^{N \times d}$ located in HBM, key block size $B_c$, where $T_c = \lceil \frac{N}{B_c} \rceil$.
      \STATE Dynamically assign registers based on the number of consumer warps.
      \STATE Allocate on-chip memory for $\mathbf{O}_i = (0) \in \mathbb{R}^{B_r \times d}$, and initialize $\ell_i, m_i = (0), (-\infty) \in \mathbb{R}^{B_r}$.
      \STATE Ensure both $\mathbf{Q}_i$ and $\mathbf{K}_0$ are present in shared memory before proceeding.
      \STATE Using WGMMA, calculate $\mathbf{S}_{\mathrm{cur}} = \mathbf{Q}_i \mathbf{K}_0^T$, then commit and synchronize.
      \STATE Free the initial stage of the buffer designated for $\mathbf{K}$.
      \STATE Derive $m_{i}$, $\tilde{\mathbf{P}}_{\mathrm{cur}}$, and $\ell_{i}$ from $\mathbf{S}_{\mathrm{cur}}$, and appropriately scale $\mathbf{O}_i$.
      \FOR{$1 \le j < T_c - 1$}
        \STATE Confirm that $\mathbf{K}_j$ has been loaded into shared memory.
        \label{code:mainloop_start}
        \STATE Compute $\mathbf{S}_{\mathrm{next}} = \mathbf{Q}_i \mathbf{K}_{j}^T$ with WGMMA; commit the result without synchronization. \label{code:first_wgmma}
        \STATE Ensure $\mathbf{V}_{j-1}$ is in shared memory.
        \STATE Update $\mathbf{O}_{i}$ by evaluating $\mathbf{O}_{i} = \mathbf{O}_{i} + \tilde{\mathbf{P}}_{\mathrm{cur}} \mathbf{V}_{j-1}$ via WGMMA; commit the output without waiting. \label{code:second_wgmma}
        \STATE Synchronize after WGMMA completes for $\mathbf{Q}_i \mathbf{K}_{j}^T$.
        \STATE Use $\mathbf{S}_{\mathrm{next}}$ to compute $m_{i}$, $\tilde{\mathbf{P}}_{\mathrm{next}}$, and $\ell_{i}$. \label{code:softmax}
        \STATE Wait until WGMMA on $\tilde{\mathbf{P}}_{\mathrm{cur}} \mathbf{V}_{j-1}$ is finished, then update the scaling of $\mathbf{O}_i$.
        \STATE Release the $(j\,\%\,s)$th buffer stage for $\mathbf{K}$ and the $(j-1\,\%\,s)$th stage for $\mathbf{V}$.
        \STATE Assign the value from $\mathbf{S}_{\mathrm{next}}$ to $\mathbf{S}_{\mathrm{cur}}$.
        \label{code:mainloop_end}
      \ENDFOR \STATE Wait until $\mathbf{V}_{T_c - 1}$ is loaded into shared memory.
      \STATE Evaluate
      $\mathbf{O}_{i} = \mathbf{O}_{i} + \tilde{\mathbf{P}}_{\mathrm{last}} \mathbf{V}_{T_c - 1}$ utilizing WGMMA, and commit while synchronizing the outcome.

\STATE Epilogue: Adjust $\mathbf{O}_{i}$ according to $m_{i}$. Determine $L_{i}$ using $m_{i}$ and $\ell_{i}$.
Store $\mathbf{O}_{i}$ and $L_{i}$ in HBM as the $i$-th segment within $\mathbf{O}$ and $L$.
\end{algorithmic}
\end{algorithm}

\cref{alg:flash3_wgmma} serves as a substitute for the consumer pathway outlined in \cref{alg:flash3_wgmma_ws_only}, thereby forming the full \textsc{FlashAttention-3} algorithm operating with FP16 accuracy. Conceptually, we employ WGMMA as shorthand for asynchronous GEMM computations. In the mainloop spanning lines \ref{code:mainloop_start} to \ref{code:mainloop_end}, the execution of the second WGMMA operation during iteration $j$ (line \ref{code:second_wgmma}) is concurrently overlapped with the softmax processing of iteration $j+1$ (line \ref{code:softmax}).

Although the pipelined design discussed above promises improved theoretical throughput, several pragmatic factors must be taken into account:

\paragraph{Compiler reordering} Although the pseudocode assumes a specific sequence of operations, compilers such as NVCC frequently reorder instructions to optimize execution. This optimization process can interrupt the intended alternation between WGMMA and non-WGMMA operations, possibly resulting in deviations from the expected flow or undermining performance improvements. Inspection of the SASS output, as detailed in Section~\ref{sec:2-stage-sass}, confirms that the compiler does in fact produce code with overlapping execution as desired.

\paragraph{Register pressure} Achieving high efficiency requires minimizing register spilling. The 2-stage pipeline architecture, however, necessitates additional registers to hold intermediary values and maintain stage information. In particular, the architecture requires storage for an additional $\mathbf{S}_{\mathrm{next}}$ within registers, incurring extra register consumption of $B_r \times B_c \times \text{sizeof}(\text{float})$ per threadblock. This elevated register requirement can hinder the use of larger block sizes—which themselves demand substantial register resources—necessitating a balance informed by empirical profiling.

\paragraph{3-stage pipelining} Building upon the previously described 2-stage approach, we introduce a 3-stage scheme that seeks to overlap the second WGMMA with the softmax computation. This further pipeline depth may yield enhanced Tensor Core utilization; however, it imposes an even greater register burden owing to the additional intermediary stage. Consequently, managing the interplay between tiling granularity and pipeline depth becomes more intricate. The complete formulation and experimental evaluation of this 3-stage strategy are presented in ~\cref{sec:3-stage}.

\subsection{Low-precision with FP8}
\label{sec:algofp8}

\textbf{Efficiency: layout transformations.} Executing the forward pass of \textsc{FlashAttention-3} using FP8 precision introduces further complexities related to layout conformity, which are not present when working with FP16.

To begin, it is important to highlight that the input tensors $\mathbf{Q}$, $\mathbf{K}$, and $\mathbf{V}$ are usually stored contiguously with respect to the head dimension. However, in order to comply with the k-major requirement of FP8 WGMMA for the second GEMM operation, $\mathbf{V}$—specifically the portions of $\mathbf{V}$ brought into SMEM—must instead be contiguous along the sequence length dimension. As TMA load operations cannot alter the contiguous axis, this necessitates either (1) pre-transposing $\mathbf{V}$ within GMEM, or (2) transposing each tile of $\mathbf{V}$ in-kernel once they are loaded into SMEM. For strategy (1), we may (1a) integrate the transpose operation into the epilogue of a prior procedure (such as rotary embedding), or (1b) invoke a dedicated pre-processing transpose kernel\footnote{Optimized transpose kernels can approach device bandwidth performance~\citep{colfax_cutlass_transpose_2024}.} to swap the strides associated with the sequence length and head dimensions. Nonetheless, (1a) presents integration challenges with standard library implementations, and (1b) is not efficient in memory-constrained inference scenarios due to excessive overhead.

In implementing FP8 \textsc{FlashAttention-3}, we instead utilize approach (2). Specifically, our in-kernel transpose leverages the LDSM (\verb|ldmatrix|) and STSM (\verb|stmatrix|) instructions, enabling warps to collaboratively transfer data between SMEM and RMEM at 128-byte segments. \footnote{The PTX documentation specifies that LDSM/STSM copy $8 \times 8$ matrices composed of 16-bit elements \cite[\S 9.7.13.4.15-16]{ptx}; to adapt these instructions for FP8, we pack two 8-bit elements into each entry, thus making these instructions compatible with FP8 data. However, since the transpose versions of LDSM/STSM cannot divide packed 8-bit elements, additional register operations are required between LDSM and STSM to achieve the correct tile-wise transpose; further technical specifics are omitted here.} These instructions maintain a low register footprint, so execution within the producer warpgroup is practical, and they support simultaneous data movement and layout transposition. Additionally, after completing the initial iteration, we can schedule the transpose operation for the upcoming $\mathbf{V}$ tile to be performed concurrently—hidden behind the computation of the two WGMMAs that operate on the preceding $\mathbf{V}$ tile and the present $\mathbf{K}$ tile.

Secondly, we note that, in contrast to FP16, the FP32 accumulator utilized in FP8 WGMMA exhibits a register memory layout that deviates from the arrangement expected for operand A in registers. Illustrative examples of these differing configurations are shown in \cref{fig:rmem_accum} and \cref{fig:rmem_operand}, with registers per thread containing entries in the indicated sequence. Through application of byte permutation instructions, it becomes possible to reformat the accumulator from the first WGMMA into an appropriate layout for subsequent WGMMA operations, aligning with the arrangement of the $\mathbf{V}$ tile generated by the in-kernel transpose. In particular, following \cref{fig:rmem_accum}, we systematically reorder registers as follows:
$$\{ \verb|d0 d1 d4 d5 d2 d3 d6 d7| \},$$
and this permutation pattern recurs for each contiguous 8-byte segment. On a structural level, this operation effectively reorders the columns of the $\mathbf{P}$ tile—for instance, the columns numbered $0189$ are repositioned as the initial four columns. To ensure the output tile from WGMMA remains accurate, we can have the in-kernel transpose write a correspondingly permuted set of rows in the $\mathbf{V}$ tile.\footnote{By conducting the transpose within the kernel, we gain increased flexibility and obviate the need for shuffle instructions to transfer register contents between threads, as previously discussed in~\citep{colfax_fp8_flashattention_2024}.}

\textbf{Accuracy: block quantization and incoherent processing.} The FP8 (e4m3) numerical format allocates 3 bits for the mantissa and 4 bits for the exponent. This minimal mantissa precision produces greater quantization errors when compared with FP16 or BF16 formats. Further complicating matters, large-scale models often exhibit outlier values~\citep{dettmers2208llm, sun2024massive}—values whose magnitude exceeds most other parameters—which complicates effective quantization. Standard approaches address this with per-tensor scaling~\citep{micikevicius2022fp8}, assigning a single scale factor to each tensor (for instance, one scale each for $\mathbf{Q}$, $\mathbf{K}$, and $\mathbf{V}$). To counteract the loss in numerical fidelity when using FP8 for attention computations, we apply the following two methods:

\begin{enumerate}

\item \textbf{Block quantization}: For block quantization, a single scalar is retained per block. Specifically, the tensors $\mathbf{Q}$, $\mathbf{K}$, and $\mathbf{V}$ are partitioned into blocks of dimensions $B_r \times d$ or $B_c \times d$, and each block undergoes quantization independently. This process can be combined with pre-attention operations, such as rotary embedding, without introducing any performance degradation, since the rotary embedding is primarily limited by memory bandwidth. Given that the \textsc{FlashAttention-3} algorithm is inherently designed to function on blocks, each block of
$\mathbf{S}$ can be rescaled accordingly to reflect the applied block quantization, incurring no extra computational overhead.

\item \textbf{Incoherent processing}: To mitigate the effects of outlier values, prior to quantizing to FP8, $\mathbf{Q}$ and $\mathbf{K}$ are multiplied by a random orthogonal matrix $\mathbf{M}$. Because $\mathbf{M}$ is orthogonal, it holds that $\mathbf{M} \mathbf{M}^\top = I$, leading to $(\mathbf{Q} \mathbf{M}) (\mathbf{K}
\mathbf{M})^\top = \mathbf{Q} \mathbf{K}^\top$. This indicates that applying $\mathbf{M}$ to both $\mathbf{Q}$ and $\mathbf{K}$ does not affect the resulting attention computation. The purpose of this transformation is to distribute outliers across the vectors, since any entry in $\mathbf{Q} \mathbf{M}$ or $\mathbf{K} \mathbf{M}$ consists of a random combination of elements from $\mathbf{Q}$ or $\mathbf{K}$, thereby diminishing quantization errors. In actual implementation, following the procedures in \citet{chee2024quip} and~\citet{tseng2024quip}, $\mathbf{M}$ is constructed as a product of random diagonal matrices whose entries are $\pm 1$ and a Hadamard matrix. This structure enables multiplication in $O(d \log d)$ complexity, as opposed to $O(d^2)$, and allows seamless integration with the rotary embedding with no extra computational burden.
\end{enumerate}

Empirical evidence in \cref{sec:numerical_error} confirms that these methods can decrease numerical error by as much as 2.6$\times$.

\section{Empirical Validation}
\label{sec:exp}

To construct and assess the efficiency and precision of \textsc{FlashAttention-3}, we employ CUTLASS~\citep{Thakkar_CUTLASS_2023} primitives, specifically leveraging abstractions like WGMMA and TMA.

\begin{itemize}

\item \textbf{Benchmarking attention.} We evaluate the execution time of \textsc{FlashAttention-3} over a range of sequence lengths and benchmark it against the baseline PyTorch implementation, \textsc{FlashAttention-2}, as well as the Triton-based version of \textsc{FlashAttention-2} (leveraging H100-specific instructions), and a cuDNN-optimized vendor release of \textsc{FlashAttention-2} for H100 GPUs. Our findings indicate that \textsc{FlashAttention-3} achieves speedups of up to 2.0$\times$ relative to \textsc{FlashAttention-2} and 1.5$\times$ compared to the Triton variant. Additionally, \textsc{FlashAttention-3} attains throughput up to 740 TFLOPs/s, reaching 75\% of H100’s peak theoretical TFLOPs/s.
\item \textbf{Ablation study.} Through systematic ablation, we demonstrate that enhancements from warp-specialization and GEMM-softmax pipelining are significant contributors to the improved efficiency of \textsc{FlashAttention-3}.
\item \textbf{Accuracy of FP8 attention.} We show that employing block quantization in conjunction with incoherent processing reduces FP8 numerical errors in \textsc{FlashAttention-3} by a factor of 2.6$\times$.

\subsection{Benchmarking Attention}
\label{subsec:benchmark_attn}

We evaluated the execution time of various attention algorithms on an H100 80GB SXM5 GPU, exploring both configurations with and without causal masking, and using head dimensions of either 64 or 128, all with FP16 input data types. The collected benchmarking results are presented in~\cref{fig:benchmark_attn_fwd} and~\cref{fig:benchmark_attn_bwd}. These figures illustrate that \textsc{FlashAttention-3} achieves approximately 1.5-2.0$\times$ faster performance than \textsc{FlashAttention-2} during the forward computation, and about 1.5-1.75$\times$ improvement in the backward computation. When compared against a conventional attention implementation, \textsc{FlashAttention-3} yields speedups of up to 3-16$\times$. Notably, for sequence lengths of 1,000 or more, \textsc{FlashAttention-3} outpaces the vendor-provided library (cuDNN, which is proprietary), even though this library is specifically tuned for H100 hardware.

\paragraph{Benchmark settings:} We vary the sequence length as 512, 1k, ..., 16k, and set batch size so that the total number of tokens is 16k. We set the hidden dimension to 2048, and head dimension to be either 64, 128, or 256 (i.e., 32 heads, 16 heads, or 8 heads). To calculate the FLOPs of the forward pass, we use:

\begin{equation*}
  4 \cdot \text{seqlen}^2 \cdot \text{head dimension} \cdot \text{number of heads}.
\end{equation*}

Employing causal masking requires halving the value to reflect that roughly 50% of the elements are computed. For the backward pass, the FLOPs are estimated by scaling the forward pass FLOPs by a factor of 2.5, given that the forward pass involves 2 matrix multiplications and the backward pass, through recomputation, involves 5 matrix multiplications.

Additionally, we assess the forward pass execution time with FP8 in comparable conditions. The outcomes for headdim 256 are presented in ~\cref{fig:benchmark_attn_fp8}, while a comprehensive set of results is provided in \cref{sec:benchmark_attn_fp8_full}.

\subsection{Ablation Study: 2-Stage Pipelining Experiments}
\label{sec:2-stage-pipelining-experiments}

We conduct ablation studies on the 2-stage WGMMA-softmax pipelining and the use of warp-specialization for the non-causal FP16 \textsc{FlashAttention-3}, maintaining fixed parameter values $\{\text{batch}, \text{seqlen}, \text{nheads}, \text{hdim}\} = \{ 4, 8448, 16, 128\}$. As illustrated in~\cref{table:ablation_pipelining}, these algorithmic enhancements—specifically, the combination of asynchronous execution with warp-specialization and the concurrent computation of GEMM and softmax—yield a substantial improvement in speed, elevating performance from 570 to 661 TFLOPs.

\subsection{Numerical Error Validation}
\label{sec:numerical_error}

Given the ongoing discussion concerning the numerical inaccuracies associated with \textsc{FlashAttention}~\citep{golden2024flash}, we evaluate \textsc{FlashAttention-2}, \textsc{FlashAttention-3}, and a conventional attention implementation by comparing each to a reference version operating in FP64 precision. To replicate the occurrence of outlier features and activations typical in LLMs~\citep{dettmers2208llm,
  sun2024massive}, we use a specific distribution to sample the elements of $\mathbf{Q}, \mathbf{K}, \mathbf{V}$ as follows:

\begin{equation*}
  \mathcal{N}(0, 1) + \mathcal{N}(0, 100) \cdot \mathrm{Bernoulli}(0.001).
\end{equation*}

In this setup, every element is sampled from a normal distribution with zero mean and unit standard deviation, except for 0.1\% of the elements, which receive an additional independent component drawn from a normal distribution with standard deviation 10. We subsequently report the root mean squared error (RMSE) in~\cref{table:numerical_error}. When using FP16 precision, both \textsc{FlashAttention-2} and \textsc{FlashAttention-3} yield RMSE values that are 1.7$\times$ lower than those observed with the conventional approach, attributed to storing intermediate values (specifically, the softmax) in FP32. The baseline implementation for FP8 leverages per-tensor scaling, conducts matrix multiplication accumulation in FP32, and maintains intermediate softmax outputs in FP16. Owing to the incorporation of block quantization and non-coherent computation, \textsc{FlashAttention-3} in FP8 achieves an accuracy improvement of 2.6$\times$ relative to this baseline.

\section{Dicussion, Limitations, Conclusion}
\label{sec:discussion}

Utilizing \textsc{FlashAttention-3}, we have shown that leveraging innovative programming methods and hardware functionalities—including asynchronous operations and low-precision arithmetic—can substantially improve both the performance and reliability of attention mechanisms. Our approach accelerates attention by 1.5–2.0$\times$ over \textsc{FlashAttention-2}, and diminishes FP8 numerical error by 2.6$\times$ compared to traditional per-tensor quantization. Nonetheless, there are several constraints in our current methodology that we intend to explore further, such as enhancing optimization for LLM inference, incorporating persistent kernel architectures in the FP8 kernel,\footnote{For benchmarking, FP16 \textsc{FlashAttention-3} employs a persistent kernel and a load balancing scheme, unlike FP8 \textsc{FlashAttention-3}, which omits these. This accounts, in part, for FP8 \textsc{FlashAttention-3}’s lower performance on tasks with small sequence lengths and causal masking, relative to FP8 cuDNN kernels.} and investigating how low-precision attention affects large-scale training. Although our experiments concentrated on Hopper GPUs, we anticipate these techniques are transferable to other types of hardware accelerators. Ultimately, we aspire for this enhanced and more precise attention primitive to catalyze advances in long-context tasks and facilitate emerging applications.

\end{document}